\chapter{Endpoint Hardening Controls}
\label{ch:4.3}
\chaptertopic{Topic 4.3}

\chapterquote{An ounce of prevention is worth a pound of cure.}{Benjamin Franklin (1736)}

\section{Motivation}

Chapter~\ref{ch:4.2} controlled who is allowed to authenticate. This chapter is about the device
itself, independent of who is logging into it. How does it arrive from the factory? What actively
watches it while it runs? What happens after a vendor discovers a flaw in it? The three sections
build a timeline: harden the device before it is ever deployed, defend it continuously while it
operates, and close the specific flaws vendors find in it over its whole operational life.

\begin{apalignment}{Unit 4: Securing Devices}{Topic 4.3 -- Protecting Devices}
\begin{apobjectives}
\apobjective{4.3.A}{Mitigate Risk}{Apply the four categories of device hardening --
changing default credentials, disabling unneeded services and ports, enabling protective features
the vendor ships off by default, and restricting accounts and privileges -- and explain how
configuration-management tooling detects and corrects the configuration drift that erodes a
hardening baseline over time.}
\apobjective{4.3.B}{Mitigate Risk}{Account for how anti-malware software actually reaches a
verdict -- by identity, by structure, or by observed behavior -- and state what each of those three
procedures can and cannot see, including why no update schedule allows a signature to match code
that has never been analyzed.}
\apobjective{4.3.C}{Mitigate Risk}{Explain why applying operating-system and application updates
reduces risk in a specific rather than general way: each fix removes the particular defect an
exploit depends on, and the interval between a public exploit and a deployed fix is the window in
which most compromise happens.}
\apobjective{4.3.D}{Mitigate Risk}{Set up a firewall on the device itself -- default action,
per-network profile, ordered or precedence-based rules restricted by port, source, and program, and
logging -- and be able to say which traffic this catches that a firewall at the network boundary
never sees.}
\end{apobjectives}
\end{apalignment}

\section{Device Hardening}

Every device ships from the factory configured for maximum compatibility and ease of setup, not
security. Identical default administrator passwords ship across thousands of units. Unneeded
services run because a manufacturer assumed someone might want them. Open remote-management ports
are left convenient for a technician.

\begin{definition}{Device hardening}
The systematic process of closing unnecessary doors before a device is deployed, and periodically
re-closing them as configurations drift. The goal is to reduce the device's \textbf{attack
surface} -- the total set of ways an attacker could interact with it to cause harm.
\end{definition}

The 2016 Mirai botnet from Chapter~\ref{ch:4.1} is the canonical illustration of what an
unhardened device costs at scale: it was built almost entirely from IoT devices still running
factory credentials nobody had ever changed.\srccite{4.3}{1} Hardening frameworks like CIS
Benchmarks and NIST SP 800-70 decompose the practice into four categories. \textbf{Change
defaults}: every default password -- the OS account, the BIOS/UEFI firmware interface, a router or
camera's built-in web panel -- must become a unique, strong credential. Default credentials are
publicly documented, so there is no secrecy to rely on. \textbf{Remove or disable what is not
needed}: every service, protocol, and open port not actively required by the device's function is
an unnecessary risk. The governing question for each one is simply whether the device's legitimate
business function requires it. \textbf{Enable protective features} that ship off by default
because they add friction: host-based firewalls, full-disk encryption, automatic screen lock,
Secure Boot to block unsigned code at startup. \textbf{Lock down accounts and privileges}: disable
built-in guest accounts, rename or disable the default administrator/root account, and require
explicit privilege elevation (\texttt{sudo}, UAC) for administrative tasks rather than running
everything as root.

\begin{lstlisting}[language=Python, caption={Checking a device configuration against a hardening baseline}]
def check_hardening(config, baseline):
    """config, baseline: dicts of the same control names -> expected
       values, e.g. {"default_password_changed": True, "telnet_enabled": False}.
       Returns sorted control names where config deviates from baseline."""
    return sorted(name for name, expected in baseline.items()
                  if config.get(name) != expected)
\end{lstlisting}

A hardened device does not stay hardened. Administrators install new software, enable a port
``temporarily'' for a project, or reset a setting during troubleshooting, and each change is a
potential re-opening of a closed door -- a phenomenon called \textbf{configuration drift}.
Professional security programs address this with configuration-management tools (Ansible, Puppet,
Group Policy) that continuously compare a device's live state against its known-good baseline and
alert on, or automatically correct, any deviation.

\begin{misconception}
\textbf{``Changing the password is enough.''} Changing the default password is the single most
important first step, but a genuinely hardened device also disables unneeded services, enables
protective features, and restricts account privileges -- password change alone leaves the other
three categories untouched.
\end{misconception}

\section{Antivirus and Endpoint Protection}

Even on a fully hardened device, a user still has to open files, visit websites, and plug in
external media -- and every one of those interactions is an opportunity for malware to run.

\begin{definition}{Signature-based detection}
A malware-detection method that matches files or memory against a database of unique byte
sequences extracted from known malware. Fast and accurate for cataloged threats, but blind by
construction to anything not yet in the database.
\end{definition}

Signature detection's blind spot is exactly what \textbf{polymorphic malware} exploits: code that
changes its shape or re-encrypts itself on every compilation, altering its signature while its
malicious function stays identical. The same threat evades a signature check indefinitely without
ever changing what it actually does. \textbf{Heuristic analysis} answers this by scanning
a program's code for patterns commonly associated with malware -- suspicious API calls, instructions
that modify critical system files -- rather than matching an exact fingerprint. \textbf{Behavioral
analysis} goes further still, watching a program \emph{as it runs}. A seemingly harmless utility
that suddenly injects code into another process, opens a raw network socket, and begins encrypting
user files in rapid succession is flagged and terminated as characteristic of ransomware -- entirely
independent of whether its file signature matches anything on record. The tradeoff is direct:
behavioral analysis catches genuinely novel threats that signatures cannot, but it is
computationally expensive and carries a higher \textbf{false-positive} rate, since legitimate
administrative software can perform some of the same actions a real attack would.

\begin{definition}{Host-based firewall}
Security software running directly on a single device -- distinct from the network firewall at an
organization's perimeter -- that filters inbound and outbound traffic against local rules.
\end{definition}

A host-based firewall closes two specific gaps a perimeter firewall cannot. First, it blocks an
already compromised machine on the same local network from reaching services on this device at
all. Second, it blocks malware that does manage to run on this device from ever reaching its
\textbf{command-and-control (C2)} server outbound. That second gap renders an otherwise-successful
infection functionally useless, because the malware can never receive instructions or exfiltrate
anything.

\begin{example}[Configuring a host-based firewall on Linux and Windows]
On a Linux endpoint, \texttt{ufw} (Uncomplicated Firewall) is a common front end for the kernel's
own \texttt{netfilter} rules. A minimal hardening baseline denies everything inbound by default and
explicitly allows only the ports the device's function actually requires:
\begin{lstlisting}[language=bash, caption={ufw baseline: deny by default, allow only SSH and HTTPS}]
sudo ufw default deny incoming
sudo ufw default allow outgoing
sudo ufw allow 22/tcp    comment 'SSH management'
sudo ufw allow 443/tcp   comment 'HTTPS service'
sudo ufw enable
\end{lstlisting}
On Windows, the equivalent baseline is set through the \texttt{NetSecurity} PowerShell module
rather than a GUI, which is what makes it scriptable across a fleet via configuration management:
\begin{lstlisting}[caption={Windows NetSecurity: block inbound by default, allow one named rule}]
Set-NetFirewallProfile -Profile Domain,Public,Private -DefaultInboundAction Block
New-NetFirewallRule -DisplayName "Allow-RDP-Mgmt" -Direction Inbound `
    -Protocol TCP -LocalPort 3389 -Action Allow -Profile Domain
\end{lstlisting}
Both examples encode the same underlying policy: a default-deny inbound posture with narrow, named
exceptions. That is precisely the traffic a perimeter firewall cannot enforce once two devices are
already talking on the same local segment.
\end{example}

\begin{lstlisting}[language=Python, caption={A detection rule engine combining signatures and behavioral thresholds}]
def evaluate_process(process, signatures, behavior_threshold=3):
    """process: dict like {"hash": str, "actions": list[str]}.
       signatures: set of known-bad hashes.
       Returns "signature-match", "behavioral-flag", or "clean"."""
    if process["hash"] in signatures:
        return "signature-match"
    suspicious_actions = {"code_injection", "raw_socket", "mass_file_encrypt"}
    matched = suspicious_actions & set(process["actions"])
    if len(matched) >= behavior_threshold:
        return "behavioral-flag"
    return "clean"
\end{lstlisting}

Managing antivirus and firewall configuration across thousands of individual devices by hand is
impractical. That is why enterprises deploy an \textbf{Endpoint Protection Platform (EPP)}: a
single agent per device combining antivirus, host firewall, intrusion prevention, and data-loss
prevention, all reporting to one central console. That console lets administrators monitor
fleet-wide status, push updates uniformly, and remotely isolate a compromised device with a single
click.

\begin{misconception}
\textbf{``Having an antivirus means I don't need a firewall.''} Antivirus blocks malicious files
and code execution; a firewall regulates network traffic. They defend against different vectors
entirely, and neither one substitutes for the other.
\end{misconception}

\section{Patching and Regular Updates}

Antivirus and firewalls defend a device from the outside, watching for known-bad files and
traffic. Many of history's most damaging breaches needed neither a tricked user nor a signature
evasion at all -- they exploited a \textbf{known, published} vulnerability that the victim simply
had not yet fixed.

\begin{definition}{CVE (Common Vulnerabilities and Exposures)}
A public, unique identifier assigned to a disclosed software vulnerability, typically accompanied
by a \textbf{CVSS} severity score from 0 to 10 describing how easy and damaging exploitation would
be.
\end{definition}

A vulnerability's lifecycle runs through discovery, responsible disclosure to the vendor with time
to fix it, and patch release. Then comes \textbf{weaponization} -- attackers reverse-engineering
the patch itself, sometimes within hours, since a patch's code diff directly reveals exactly what
was broken and how. The moment a patch ships, the vulnerability effectively becomes public
knowledge regardless of whether an exploit yet exists. A race begins between defenders applying
the fix and attackers building a working exploit from the disclosure.

\begin{figure}[h]
\centering
\begin{tikzpicture}[
  ev/.style={draw=kickstartermaroon, thick, rounded corners, minimum width=2.4cm, minimum height=1.0cm, align=center, font=\scriptsize},
]
\node[ev, fill=kickstartercream] (disc) at (0,0) {Discovery};
\node[ev, fill=kickstartercream] (disclose) at (2.9,0) {Responsible\\disclosure};
\node[ev] (patch) at (5.8,0) {Patch\\released};
\node[ev, fill=kickstartermaroon!20] (weapon) at (8.7,0) {Weaponization\\(diff reverse-\\engineered)};
\node[ev, fill=kickstartermaroon!30] (exploit) at (11.6,0) {n-day exploit\\against unpatched\\systems};
\draw[-{Latex[length=2mm]}, kickstartermaroon] (disc) -- (disclose);
\draw[-{Latex[length=2mm]}, kickstartermaroon] (disclose) -- (patch);
\draw[-{Latex[length=2mm]}, kickstartermaroon] (patch) -- (weapon);
\draw[-{Latex[length=2mm]}, kickstartermaroon] (weapon) -- (exploit);
\node[font=\scriptsize\sffamily, gray] at (5.8,-1.1) {a zero-day attack occurs before this point; every stage after it is the n-day race};
\end{tikzpicture}
\caption{The vulnerability lifecycle. Patch release does not end the risk -- it starts a race
between defenders applying the fix and attackers reverse-engineering it from the patch itself.}
\label{fig:4.3-vuln-lifecycle}
\end{figure}

\begin{example}[Zero-day versus n-day]
A \textbf{zero-day} vulnerability is being actively exploited \emph{before} any patch exists at
all. There are zero days of protection available even for a fully up-to-date system, and only
other layered defenses (behavioral detection, network segmentation) can help. An \textbf{n-day}
vulnerability is one for which a patch has existed for some number of days, but the target system
simply has not applied it yet. The overwhelming majority of real-world breaches exploit n-day
vulnerabilities, not zero-days. Reusing a known exploit against an unpatched target is far cheaper
for an attacker than discovering something genuinely new, and most organizations lag behind on
patching long enough to make that bet reliably pay off.
\end{example}

\begin{example}[CVE-2017-0144 and the WannaCry exploit window]
\textbf{CVE-2017-0144}, an SMBv1 remote-code-execution flaw in Windows (CVSS v3 base score 8.8),
illustrates the patch-to-weaponization race precisely. Microsoft shipped the fix as bulletin
\textbf{MS17-010} on \textbf{March 14, 2017}.\srccite{4.3}{2} The exploit built around that flaw,
code-named \textbf{EternalBlue}, was leaked publicly by the Shadow Brokers group on April 14,
2017 -- one month after the patch. Just under two months after the original patch, on
\textbf{May 12, 2017}, the \textbf{WannaCry} ransomware worm combined EternalBlue with a
self-propagation mechanism to infect more than 200{,}000 machines across 150
countries within days.\srccite{4.3}{3} Every one of those machines had gone unpatched despite the
fix already sitting on Microsoft's servers for that entire window. WannaCry is not a
zero-day case study; every infected machine was vulnerable to a flaw with a known fix, which is
exactly what makes it the canonical n-day breach rather than a genuinely novel attack.
\end{example}

Patching is not free, which is the actual reason unpatched year-old vulnerabilities remain common
despite patching's obvious value. A patch can break a business-critical application that depends
on old behavior. Some patches require a reboot that must be scheduled around business needs. A
large enterprise may have tens of thousands of endpoints, each needing separate validation. Some
critical systems run software the vendor no longer supports at all (\textbf{end-of-life}), meaning
no patch will ever exist and only compensating controls like network isolation remain available.

\begin{lstlisting}[language=Python, caption={Prioritizing a patch by severity, exploitation status, and exposure}]
def patch_urgency(cvss_score, actively_exploited, internet_facing):
    """Returns "emergency", "high", "scheduled", or "routine"."""
    if actively_exploited and internet_facing:
        return "emergency"
    if actively_exploited or (cvss_score >= 9.0 and internet_facing):
        return "high"
    if cvss_score >= 7.0:
        return "scheduled"
    return "routine"
\end{lstlisting}

Mature organizations run a formal \textbf{patch management} process rather than patching ad hoc.
First, inventory what is running, since a patch cannot be applied to an unknown asset. Next,
prioritize new patches by CVSS severity, whether exploitation is already occurring in the wild,
and whether the affected service is internet-facing. Test in a non-production environment to catch
compatibility breaks before they hit production. Deploy in waves on a defined \textbf{patch
cadence}, reserving emergency out-of-cycle deployment for critical, actively exploited flaws.
Finally, verify the patch was actually applied fleet-wide -- partial or failed rollouts are common,
and they leave silent, undetected gaps behind.

\begin{misconception}
\textbf{``Zero-days are the main way organizations get breached.''} Most real-world breaches exploit
n-day vulnerabilities -- known flaws with an available patch that simply was not applied in time.
Zero-days are comparatively rare and expensive for an attacker to obtain or develop.
\end{misconception}

\section*{Review Questions}
\begin{reviewquestions}
\item A newly purchased IoT camera has its default administrator password unchanged and Telnet
enabled by default. Using the four-category hardening checklist from this chapter, identify which
two categories this failure falls under and explain why changing only the password would leave the
device still exposed.
\item Explain why polymorphic malware defeats signature-based detection while leaving its actual
malicious behavior completely unchanged, and explain which detection method in this chapter is
specifically designed to catch it anyway.
\item Malware successfully executes on a hardened, antivirus-protected laptop but never manages to
contact its command-and-control server. Explain which specific control is responsible for that
outcome and why it renders the infection functionally useless despite having already run.
\item A CVE is disclosed alongside its patch. Explain why attackers can sometimes weaponize the
vulnerability within hours of the patch's release, even though the patch was meant to fix the
problem, not enable new attacks.
\item Explain why an organization might rationally choose to delay applying a released patch for
several days rather than deploying it immediately, and explain how a formal patch-management
process reduces the risk of that delay without eliminating it entirely.
\item Using the vulnerability-lifecycle diagram, explain why WannaCry is placed at the
``n-day exploit'' stage rather than the ``weaponization'' stage, even though EternalBlue itself
was built by reverse-engineering a patch.
\item A CISO argues that since most breaches are n-day rather than zero-day, patching quickly
matters more than funding zero-day threat research. Using the vulnerability-lifecycle diagram and
the patch-urgency logic from this chapter, explain what evidence in this chapter supports that
argument, and identify one category of system (from the patching-is-not-free discussion) where
speed alone cannot close the gap.
\end{reviewquestions}

\begin{answerkey}
\item The unchanged default password falls under \emph{change defaults}, and the enabled Telnet
service falls under \emph{remove or disable what is not needed}. Changing only the password would
leave Telnet -- an insecure, plaintext remote-management protocol -- still reachable, giving an
attacker a second independent path onto the device that the password change never touched.
\answerlabel{4.3.A}
\item Polymorphic malware re-encrypts or reshapes its own code on every compilation, changing its
byte-level signature while its underlying malicious function stays identical, so a signature
database built on the old fingerprint never matches the new one. Behavioral analysis is designed to
catch it anyway, since it watches what a program actually does at runtime -- injecting code,
opening raw sockets, mass-encrypting files -- rather than matching against a static fingerprint.
\answerlabel{4.3.B}
\item A host-based firewall is responsible: it blocks outbound traffic from the infected device to
any command-and-control server the malware attempts to reach. This renders the infection
functionally useless despite having already executed, because the malware can neither receive
further instructions nor exfiltrate any data it collected.
\answerlabel{4.3.D}
\item A released patch's binary diff directly reveals which code paths changed. That lets an
attacker reverse-engineer exactly what was broken and reconstruct a working exploit from the
comparison alone -- as CVE-2017-0144 showed, where EternalBlue was built from the same
vulnerability class MS17-010 had already patched. The patch fixes the flaw for defenders who apply
it, but its public release simultaneously hands attackers a precise roadmap to the flaw for anyone
who has not.
\answerlabel{4.3.C}
\item A patch can break a business-critical application depending on old behavior, may require a
disruptive reboot, and at enterprise scale needs individual validation across tens of thousands of
endpoints. Immediate blanket deployment carries its own operational risk. A formal patch
management process reduces, without eliminating, that delay's risk: it prioritizes the patch using
CVSS severity and active-exploitation status, tests it in a non-production environment first, and
deploys it on a defined cadence, with emergency out-of-cycle deployment reserved for critical,
actively exploited flaws.
\answerlabel{4.3.C}
\item Weaponization is the act of reverse-engineering a working exploit from the patch, which
happened quickly for EternalBlue. But WannaCry itself did not strike until May 12, 2017, nearly two
months after the March 14 patch, and every machine it infected was one that had not yet applied a
fix that had been available the entire time. WannaCry is the exploitation of that already-weaponized
flaw against unpatched targets, which is precisely the definition of the n-day stage: the exploit
existed and worked, but the window it succeeded in was created by delayed patching, not by any
remaining novelty in the vulnerability itself.
\answerlabel{4.3.C}
\item The chapter's own claim that most real-world breaches exploit n-day rather than zero-day
vulnerabilities supports the CISO's argument, as does the vulnerability-lifecycle diagram: it shows
the n-day window opening the moment a patch ships and staying open until every affected system is
patched. Closing that window faster removes far more real-world risk than defending against the
comparatively rare zero-day case. The category where speed alone cannot close the gap is
end-of-life systems. A vendor that no longer supports the software will never release a patch to
deploy quickly in the first place, so no patch-management process, however fast, can close that
specific n-day window -- only compensating controls like network isolation remain available.
\answerlabel{4.3.C}
\end{answerkey}

\section*{Sources}
\begin{enumerate}[label=\lbrack\arabic*\rbrack, leftmargin=2.2em]
\item \srcitem{4.3}{1}M.\ Antonakakis, T.\ April, M.\ Bailey, et al., ``Understanding the Mirai
Botnet,'' \emph{26th USENIX Security Symposium}, Vancouver, August 16--18, 2017, pp.\ 1093--1110 --
the peer-reviewed analysis documenting Mirai's compromise of IoT devices through roughly 60
hardcoded factory-default Telnet username/password pairs.
\item \srcitem{4.3}{2}Microsoft Corporation, Security Bulletin MS17-010, ``Security Update for
Microsoft Windows SMB Server,'' released March 14, 2017, addressing CVE-2017-0144 (CVSS v3 base
score 8.8).
\item \srcitem{4.3}{3}Cloudflare, ``What Was the WannaCry Ransomware Attack?,'' Cloudflare Learning
Center -- reporting the Shadow Brokers' April 14, 2017 public leak of the EternalBlue exploit and
WannaCry's spread to more than 200{,}000 computers across over 150 countries beginning May 12,
2017.
\end{enumerate}
